\section{Guía de implementación}

\subsection{Configuración recomendada}

Basado en las discusiones del conferenciante y la experiencia práctica con DPM-Solver, a continuación se presentan las configuraciones recomendadas para diferentes escenarios:

\begin{table}[H]
\centering
\begin{tabular}{lcc}
\toprule
\textbf{Escenario} & \textbf{Método recomendado} & \textbf{Pasos} \\
\midrule
Vista rápida / Aplicación interactiva & DPM-Solver++(2M) & 10$\sim$15 \\
Generación de alta calidad & DPM-Solver++(2M) & 20$\sim$25 \\
Máxima calidad (sin considerar velocidad) & DPM-Solver++(3M) & 25$\sim$50 \\
\bottomrule
\end{tabular}
\caption{Configuraciones recomendadas de DPM-Solver para diferentes escenarios (2M = método multietapa de segundo orden, 3M = método multietapa de tercer orden)}
\end{table}

\subsection{Puntos clave de implementación}

\begin{warningbox}{Trampas comunes durante la implementación}
\begin{enumerate}
    \item \textbf{Estabilidad numérica}: $e^{\lambda}$ puede desbordarse cuando los valores de $\lambda$ son grandes; es necesario utilizar el cálculo en log-space.
    \item \textbf{Dirección del paso temporal}: En los modelos de difusión, el proceso "hacia adelante" implica añadir ruido ($t: 0 \to T$) y el "hacia atrás", eliminarlo ($t: T \to 0$). La dirección de resolución de la EDO es hacia atrás; $h = \lambda_t - \lambda_s$ suele ser un valor\textbf{ positivo} (ya que $\lambda$ aumenta a medida que se elimina el ruido).
    \item \textbf{Precalentamiento del método multietapa}: El método multietapa de segundo orden no tiene historial en cache durante el primer paso, por lo que debe retroceder al método unietapa de primer orden.
    \item \textbf{Compatibilidad con Classifier-free Guidance}: La guía sin clasificador altera la EDO efectiva; DPM-Solver++ cuenta con una adaptación específica para esto.
\end{enumerate}
\end{warningbox}

\begin{knowledgebox}{Integración de DPM-Solver en frameworks principales}
DPM-Solver y DPM-Solver++ han sido ampliamente integrados en las bibliotecas principales de modelos de difusión:
\begin{itemize}
    \item \textbf{Diffusers} (Hugging Face): \texttt{DPMSolverMultistepScheduler}
    \item \textbf{k-diffusion}: como uno de los muestreadores por defecto.
    \item \textbf{Stable Diffusion WebUI}: disponible en el menú desplegable de Sampler.
\end{itemize}
En la práctica, generalmente solo es necesario cambiar el scheduler/muestreador a DPM-Solver++ para disfrutar del efecto de aceleración, sin necesidad de modificar el modelo en sí.
\end{knowledgebox}

\subsection{Resumen de la sección}

El método multietapa de segundo orden (2M) de DPM-Solver++ es la opción predeterminada óptima para aplicaciones prácticas, equilibrando velocidad y calidad. Al implementarlo, se deben prestar atención a detalles como la estabilidad numérica y el precalentamiento del método multietapa.


